% Source: Wald GR, physical PDF pp. 175--181
%         (printed pp. 162--168).
% Coverage: Section 7.1, Stationary, Axisymmetric Solutions;
%           theorem 7.1.1 and equations (7.1.1)--(7.1.29).
% Figures/tables/footnotes: footnote 1; no figures or tables.
% Status: source-reviewed and compile-clean.
% Uncertainties: none.

\subsection{Stationary, Axisymmetric Solutions}\label{sec:7.1}

As already mentioned in section 6.1, a spacetime is said to be \emph{stationary} if
there exists a one-parameter group of isometries \(\sigma_t\) whose orbits are
timelike curves. Thus, every stationary spacetime possesses a timelike Killing vector
field \(\xi^a\) (see appendix C). (Conversely, every spacetime with a timelike
Killing vector field whose orbits are complete is stationary.) Similarly, we call a
spacetime \emph{axisymmetric} if there exists a one-parameter group of isometries
\(\chi_\phi\) whose orbits are closed spacelike curves, which implies the existence
of a spacelike Killing field \(\psi^a\) whose integral curves are closed. We call a
spacetime stationary and axisymmetric if it possesses both these symmetries and if,
in addition, the actions of \(\sigma_t\) and \(\chi_\phi\) commute:
\begin{equation}
  \sigma_t\circ\chi_\phi=\chi_\phi\circ\sigma_t ,
  \tag{7.1.1}\label{eq:7.1.1}
\end{equation}
i.e., the rotations commute with the time translations. This is easily seen to be
equivalent to the condition that the Killing vector fields \(\xi^a\) and \(\psi^a\)
commute,
\begin{equation}
  [\xi,\psi]=0 .
  \tag{7.1.2}\label{eq:7.1.2}
\end{equation}

Stationary, axisymmetric spacetimes are of considerable interest in general
relativity since they describe equilibrium configurations of axisymmetric, rotating
bodies. Thus, in particular, in order to generalize the results of section 6.2 to
account for rotation, we must be able to solve Einstein's equation with perfect fluid
source in the presence of these symmetries. Stationary, axisymmetric vacuum
solutions are also of great interest, since they describe the exterior gravitational
field of rotating bodies and would be needed to match onto the interior solutions.

The commutativity of \(\xi^a\) and \(\psi^a\) implies that we can choose coordinates
\((x^0=t,x^1=\phi,x^2,x^3)\) so that both
\(\xi^a=(\partial/\partial t)^a\) and
\(\psi^a=(\partial/\partial\phi)^a\) are coordinate vector fields. As discussed in
appendix C, the metric components in such a coordinate system will be independent
of \(t\) and \(\phi\), so the metric will take the form
\begin{equation}
  \dd s^2=\sum_{\mu,\nu}g_{\mu\nu}(x^2,x^3)\dd x^\mu\dd x^\nu .
  \tag{7.1.3}\label{eq:7.1.3}
\end{equation}
Thus, we must solve for 10 unknown functions, \(g_{\mu\nu}\), of two variables. We
shall show, now, that by a further careful selection of coordinate system, a weak
further assumption (see hypothesis (i) of theorem 7.1.1 below) and some use of
Einstein's equation, we can reduce the metric to a form involving only three
functions of two variables (see equation~\eqref{eq:7.1.22} below), and reduce the
problem of solving Einstein's equation to that of solving two equations for two
unknown functions (see equations~\eqref{eq:7.1.24} and \eqref{eq:7.1.25} below)
together with a quadrature for a third function (see equations~\eqref{eq:7.1.26} and
\eqref{eq:7.1.27}).

The first crucial simplification arises from the fact that under the hypotheses of the
following theorem, the two-dimensional subspaces of the tangent space at each point
which are spanned by the vectors orthogonal to \(\xi^a\) and \(\psi^a\) are
integrable, i.e., tangent to two-dimensional surfaces. We shall state and prove the
theorem first and then discuss its implications for simplifying the general form of
stationary, axisymmetric metrics.

\medskip
\noindent\textsc{Theorem 7.1.1.}\label{thm:7.1.1}\quad
\emph{Let \(\xi^a\) and \(\psi^a\) be two commuting Killing fields such that
(i) \(\xi_{[a}\psi_b\nabla_c\xi_{d]}\) and
\(\xi_{[a}\psi_b\nabla_c\psi_{d]}\) each vanishes at at least one point of the
spacetime (which, in particular, will be true if either \(\xi^a\) or \(\psi^a\)
vanishes at one point) and
(ii) \(\xi^aR_a{}^{[b}\xi^c\psi^{d]}=\psi^aR_a{}^{[b}\xi^c\psi^{d]}=0\).
Then the 2-planes orthogonal to \(\xi^a\) and \(\psi^a\) are integrable.}

\begin{proof}
This theorem is a direct application of Frobenius's theorem (see appendix B). We
want to know when the planes orthogonal to the 1-forms \(\xi_a\) and \(\psi_a\) are
integrable. According to the dual formulation of Frobenius's theorem (theorem
B.3.2) the necessary and sufficient conditions for integrability are
\begin{align}
  \nabla_{[a}\xi_{b]}&=(\mu^1)_{[a}\xi_{b]}+(\mu^2)_{[a}\psi_{b]} ,
    \tag{7.1.4}\label{eq:7.1.4}\\
  \nabla_{[a}\psi_{b]}&=(\sigma^1)_{[a}\xi_{b]}+(\sigma^2)_{[a}\psi_{b]} ,
    \tag{7.1.5}\label{eq:7.1.5}
\end{align}
where \((\mu^1)_a,(\mu^2)_a,(\sigma^1)_a,(\sigma^2)_a\) are arbitrary 1-forms. We
shall establish the theorem by proving that equation~\eqref{eq:7.1.4} holds; the proof
that equation~\eqref{eq:7.1.5} also holds follows in an identical manner.
Equation~\eqref{eq:7.1.4} is equivalent to
\begin{equation}
  \psi_{[a}\xi_b\nabla_c\xi_{d]}=0 ,
  \tag{7.1.6}\label{eq:7.1.6}
\end{equation}
which, in turn, is equivalent to
\begin{equation}
  \epsilon_{abcd}\psi^a\xi^b\nabla^c\xi^d=0 ,
  \tag{7.1.7}\label{eq:7.1.7}
\end{equation}
where \(\epsilon_{abcd}\) is the volume element associated with the metric (see
appendix B). We define the \emph{twist}, \(\omega_a\), of \(\xi^a\) by
\begin{equation}
  \omega_a=\epsilon_{abcd}\xi^b\nabla^c\xi^d .
  \tag{7.1.8}\label{eq:7.1.8}
\end{equation}
(According to equation~\eqref{eq:B.3.6}, \(\omega_a\) measures the failure of
\(\xi^a\) to be hypersurface orthogonal.) Our task is to show that
\(\psi^a\omega_a=0\).

By hypothesis (i), we know \(\psi^a\omega_a\) vanishes at at least one point, so it
will vanish everywhere if and only if its derivative is identically zero. We have
\begin{align}
  \nabla_b(\psi^a\omega_a)
    &=\psi^a\nabla_b\omega_a+\omega_a\nabla_b\psi^a\notag\\
    &=\psi^a\nabla_a\omega_b+\omega_a\nabla_b\psi^a
      +2\psi^a\nabla_{[b}\omega_{a]}\notag\\
    &=\mathcal L_\psi\omega_b+2\psi^a\nabla_{[b}\omega_{a]} ,
  \tag{7.1.9}\label{eq:7.1.9}
\end{align}
where the formula~\eqref{eq:C.2.13} for the Lie derivative was used. However, the
group of diffeomorphisms \(\chi_\phi\) generated by \(\psi^a\) leaves \(\xi^a\)
invariant (since \(\psi^a\) and \(\xi^a\) commute) and leaves the metric invariant
(since \(\psi^a\) is a Killing field). Hence, \(\chi_\phi\) must leave invariant any
tensor field that can be constructed uniquely out of \(\xi^a\) and the metric. Since
\(\omega_b\) is such a tensor field, we must have\footnote{Another way of seeing
this is that in a coordinate system adapted to both \(\xi^a\) and \(\psi^a\), any
tensor \(T^{a\dots b}{}_{c\dots d}\) whose components
\(T^{\alpha\dots\beta}{}_{\gamma\dots\delta}\) are constructed out of the
components \(\xi^\alpha\) and \(g_{\mu\nu}\), and their coordinate derivatives,
must satisfy
\(\partial(T^{\alpha\dots\beta}{}_{\gamma\dots\delta})/\partial\phi=0\), which
proves that \(\mathcal L_\psi T^{a\dots b}{}_{c\dots d}=0\) according to
equation~\eqref{eq:C.2.4}. Note that \(\mathcal L_\psi\omega_b=0\) holds, of
course, also could be proven directly using the formulas of appendix C (problem 1).}
\begin{equation}
  \mathcal L_\psi\omega_b=0 .
  \tag{7.1.10}\label{eq:7.1.10}
\end{equation}

Thus, to complete the proof, we need only compute \(\nabla_{[a}\omega_{b]}\). We
have
\begin{align}
  \epsilon^{abcd}\nabla_c\omega_d
    &=\epsilon^{abcd}\epsilon_{defg}\nabla_c(\xi^e\nabla^f\xi^g)\notag\\
    &=6\nabla_c(\xi^{[c}\nabla^a\xi^{b]}) ,
  \tag{7.1.11}\label{eq:7.1.11}
\end{align}
where equations~\eqref{eq:B.2.11} and \eqref{eq:B.2.13} were used. Now, using
Killing's equation~\eqref{eq:C.3.1}, we have
\begin{equation}
  \xi^{[c}\nabla^a\xi^{b]}=\frac13
  \{\xi^c\nabla^a\xi^b+\xi^b\nabla^c\xi^a+\xi^a\nabla^b\xi^c\} .
  \tag{7.1.12}\label{eq:7.1.12}
\end{equation}
However, we find
\begin{align}
  \nabla_c(\xi^c\nabla^a\xi^b)
    &=(\nabla_c\xi^c)\nabla^a\xi^b+\xi^c\nabla_c\nabla^a\xi^b\notag\\
    &=-\xi^cR^{ab}{}_{cd}\xi^d\notag\\
    &=0 .
  \tag{7.1.13}\label{eq:7.1.13}
\end{align}
Here the trace of Killing's equation was used to eliminate the first term, while
equation~\eqref{eq:C.3.6} was used to express the second term in terms of the
Riemann tensor. The antisymmetry of the Riemann tensor in its last two indices was
used to get the final conclusion. On the other hand, we have
\begin{align}
  \nabla_c(\xi^b\nabla^c\xi^a+\xi^a\nabla^b\xi^c)
    &=(\nabla_c\xi^b)(\nabla^c\xi^a)
      +(\nabla_c\xi^a)(\nabla^b\xi^c)
      +\xi^b\nabla_c\nabla^c\xi^a
      +\xi^a\nabla_c\nabla^b\xi^c\notag\\
    &=2\xi^{[b}\nabla_c\nabla^{|c|}\xi^{a]}\notag\\
    &=-2\xi^{[b}R^{a]}{}_c\xi^c ,
  \tag{7.1.14}\label{eq:7.1.14}
\end{align}
where Killing's equation was used to get the second line and
equation~\eqref{eq:C.3.9} was used to get the third line. Thus, we find
\begin{align}
  \nabla_{[a}\omega_{b]}
    &=-\frac14\epsilon_{abcd}\epsilon^{cdef}\nabla_e\omega_f\notag\\
    &=-\epsilon_{abcd}\xi^cR^d{}_e\xi^e ,
  \tag{7.1.15}\label{eq:7.1.15}
\end{align}
and, finally,
\begin{align}
  \nabla_b(\psi^a\omega_a)
    &=2\psi^a\nabla_{[b}\omega_{a]}\notag\\
    &=-2\epsilon_{bacd}\psi^a\xi^cR^d{}_e\xi^e\notag\\
    &=0
  \tag{7.1.16}\label{eq:7.1.16}
\end{align}
by hypothesis (ii). This completes the proof.
\end{proof}

The hypotheses of theorem 7.1.1 will be satisfied in a wide range of stationary,
axisymmetric spacetimes of physical interest. In particular, if the spacetime is
asymptotically flat, there must be a ``rotation axis'' on which \(\psi^a\) vanishes,
so hypothesis (i) will be satisfied. For a vacuum spacetime, we have \(R_{ab}=0\),
so hypothesis (ii) is trivially satisfied. In addition, hypothesis (ii) will also be
satisfied when the matter stress-energy
\(T_{ab}=(8\pi)^{-1}(R_{ab}-\tfrac12g_{ab}R)\) has the form of a perfect fluid
with 4-velocity in the plane spanned by \(\xi^a\) and \(\psi^a\) (i.e., circular
flow) or if \(T_{ab}\) is the stress tensor of a stationary, axisymmetric
electromagnetic field (Carter 1969).

For spacetimes which satisfy the hypotheses of theorem 7.1.1, we may choose
coordinates \(x^2,x^3\) in one of the orthogonal two-surfaces and ``carry'' these
coordinates to the rest of the spacetime along the integral curves of \(\xi^a\) and
\(\psi^a\). In the coordinates \((t,\phi,x^2,x^3)\), the metric components take the
form
\begin{equation}
  g_{\mu\nu}=
  \begin{pmatrix}
    -V&W&0&0\\
    W&X&0&0\\
    0&0&g_{22}&g_{23}\\
    0&0&g_{23}&g_{33}
  \end{pmatrix} ,
  \tag{7.1.17}\label{eq:7.1.17}
\end{equation}
where \(V=-g_{00}=-\xi_a\xi^a\), \(W=g_{01}=\xi^a\psi_a\),
\(X=g_{11}=\psi^a\psi_a\), and the \(2\times2\) block of zeros expresses the
orthogonality of \(\partial/\partial x^2\) and \(\partial/\partial x^3\) with
\(\partial/\partial t\) and \(\partial/\partial\phi\). (A physical interpretation of
\(W\) is established in problem 3.) Thus, theorem 7.1.1 allows us to reduce the
number of nonvanishing metric components to six. Without theorem 7.1.1, only two
metric components in the \(2\times2\) block could be set to zero by use of the
coordinate freedom available in a coordinate system adapted to \(\xi^a\) and
\(\psi^a\).

We still have not specified how the two-surface coordinates \(x^2\) and \(x^3\) are
to be chosen, and, as we shall see, significant further simplification can be achieved
by a judicious choice of these coordinates. We define the scalar function \(\rho\) by
\begin{equation}
  \rho^2=VX+W^2 ;
  \tag{7.1.18}\label{eq:7.1.18}
\end{equation}
i.e., \(\rho^2\) is minus the determinant of the \(t\)-\(\phi\) part of the metric.
Assuming \(\nabla_a\rho\ne0\), we choose \(\rho\) as one of the coordinates,
\(x^2\), of the two-surface. We choose the other coordinate, \(z=x^3\), so that
\(\nabla_az\) is orthogonal to \(\nabla_a\rho\). [This is accomplished by setting
\(z=\text{constant}\) along the integral curves of \(\nabla^a\rho\), which uniquely
determines \(z\) up to the transformations \(z\longrightarrow z'=f(z)\).] In the
coordinates \(t,\phi,\rho,z\) the metric takes the form
\begin{equation}
  \dd s^2=-V(\dd t-w\dd\phi)^2+V^{-1}\rho^2\dd\phi^2
    +\Omega^2(\dd\rho^2+\Lambda\dd z^2) ,
  \tag{7.1.19}\label{eq:7.1.19}
\end{equation}
where \(w=W/V\). Thus, we have reduced the unknown metric components to four
functions, \(V,w,\Omega,\Lambda\) of two variables \(\rho,z\).
Equation~\eqref{eq:7.1.19} is the general form of a stationary, axisymmetric
spacetime satisfying the hypotheses of theorem 7.1.1.

The form of the metric can be simplified further for vacuum spacetimes,
\(R_{ab}=0\). The components of \(R_{ab}\) in the plane spanned by \(\xi^a\) and
\(\psi^a\) can be computed from equation~\eqref{eq:C.3.9} or, alternatively, from
the general coordinate basis formulas of section 3.4a. The equation
\begin{equation}
  0=R^t{}_t+R^\phi{}_\phi
   =(\nabla_at)R^a{}_b\xi^b+(\nabla_a\phi)R^a{}_b\psi^b
  \tag{7.1.20}\label{eq:7.1.20}
\end{equation}
yields
\begin{equation}
  D^aD_a\rho=0 ,
  \tag{7.1.21}\label{eq:7.1.21}
\end{equation}
where \(D_a\) is the covariant derivative operator on the two-dimensional surface
spanned by \(\rho\) and \(z\) with the induced metric
\(\dd s^2=\Omega^2(\dd\rho^2+\Lambda\dd z^2)\). Thus, \(\rho\) is harmonic,
i.e., it satisfies the two-dimensional Laplace equation in the two-surfaces. This has
two important consequences: (1) If \(\rho\ne\text{constant}\), it can be shown that
\(\nabla_a\rho\) can vanish only at isolated points. Since our coordinates \(\rho,z\)
will be well behaved except where \(\nabla_a\rho=0\), this shows that our coordinate
system can break down only at isolated points. In fact, in many situations it is
possible to show that \(\nabla_a\rho\ne0\) everywhere, so the coordinates \(\rho,z\)
are globally well behaved (Carter 1973). (2) It follows directly from
equation~\eqref{eq:7.1.21} that \(\Lambda\) is a function of \(z\) alone. Thus we may
use the remaining coordinate freedom to transform \(z\) via
\(z\longrightarrow\int\Lambda^{1/2}\dd z\) and thereby set \(\Lambda=1\). (Note
that \(z\) then becomes the harmonic function conjugate to \(\rho\); see problem 2
of chapter 3.) Except for the trivial constant rescaling or shift of origin of the
coordinates \(t,\phi,z\), we have now completely specified our coordinate system.
In these coordinates the metric takes the remarkably simple form (Papapetrou 1953,
1966)
\begin{equation}
  \dd s^2=-V(\dd t-w\dd\phi)^2
    +V^{-1}[\rho^2\dd\phi^2+e^{2\gamma}(\dd\rho^2+\dd z^2)] ,
  \tag{7.1.22}\label{eq:7.1.22}
\end{equation}
where \(\gamma=\tfrac12\ln(V\Omega^2)\). Note that in flat spacetime
(\(V=1,w=\gamma=0\)), the coordinates \((\phi,\rho,z)\) are ordinary cylindrical
coordinates.

In deriving equation~\eqref{eq:7.1.22}, we have used Einstein's equation thus far
only in hypothesis (ii) of theorem 7.1.1 and in equation~\eqref{eq:7.1.20}. The
remaining components of the vacuum Einstein's equation \(R_{ab}=0\) can be
computed in a straightforward manner using the methods of section 3.4a. We shall
not present the details here but will merely quote the final results. We obtain four
independent equations for \(V,w\), and \(\gamma\). The first two involve only
\(V\) and \(w\) and are most conveniently formulated by defining an (unphysical)
flat three-dimensional metric,
\begin{equation}
  \dd s^2=\rho^2\dd\phi^2+\dd\rho^2+\dd z^2 ,
  \tag{7.1.23}\label{eq:7.1.23}
\end{equation}
and expressing the equations in terms of the flat derivative operator \(\boldsymbol{\nabla}\)
associated with this metric. (We use the usual vector notation \(\boldsymbol{\nabla}\) rather
than the index notation to avoid confusion with the derivative operator associated
with the true spacetime metric.) In this way, the first two equations may be viewed
as equations for axisymmetric scalar fields \(V,w\) in the unphysical
three-dimensional flat space of equation~\eqref{eq:7.1.23}. The equations are:
\begin{align}
  \boldsymbol{\nabla}\mathbin{\cdot}
    \{V^{-1}\boldsymbol{\nabla} V+\rho^{-2}V^2w\boldsymbol{\nabla} w\}&=0 ,
    \tag{7.1.24}\label{eq:7.1.24}\\
  \boldsymbol{\nabla}\mathbin{\cdot}
    \{\rho^{-2}V^2\boldsymbol{\nabla} w\}&=0 ,
    \tag{7.1.25}\label{eq:7.1.25}\\
  \frac{\partial\gamma}{\partial\rho}
    &=\frac14\rho V^{-2}
      \left[\left(\frac{\partial V}{\partial\rho}\right)^2
            -\left(\frac{\partial V}{\partial z}\right)^2\right]\notag\\
    &\quad-\frac14\rho^{-1}V^2
      \left[\left(\frac{\partial w}{\partial\rho}\right)^2
            -\left(\frac{\partial w}{\partial z}\right)^2\right] ,
    \tag{7.1.26}\label{eq:7.1.26}\\
  \frac{\partial\gamma}{\partial z}
    &=\frac12\rho V^{-2}\frac{\partial V}{\partial\rho}
      \frac{\partial V}{\partial z}
      -\frac12\rho^{-1}V^2\frac{\partial w}{\partial\rho}
      \frac{\partial w}{\partial z} .
    \tag{7.1.27}\label{eq:7.1.27}
\end{align}

The last two equations are in danger of overdetermining \(\gamma\). However, the
integrability condition
\(\partial^2\gamma/\partial z\partial\rho=
\partial^2\gamma/\partial\rho\partial z\) is identically satisfied by virtue of
equations~\eqref{eq:7.1.24} and \eqref{eq:7.1.25}. Hence, given a solution of
equations~\eqref{eq:7.1.24} and \eqref{eq:7.1.25}, a solution of
equations~\eqref{eq:7.1.26} and \eqref{eq:7.1.27} always exists and is unique up to
the addition of a constant. This solution can be obtained by performing a line
integral of the right sides of equations~\eqref{eq:7.1.26} and \eqref{eq:7.1.27},
i.e., explicitly,
\begin{equation}
  \gamma(p)-\gamma(q)=\int_q^p
    \left(\frac{\partial\gamma}{\partial\rho}\dd\rho
          +\frac{\partial\gamma}{\partial z}\dd z\right) .
  \tag{7.1.28}\label{eq:7.1.28}
\end{equation}

Thus, apart from the computation required to find \(\gamma\) explicitly, the problem
of solving for all stationary, axisymmetric vacuum solutions of Einstein's equation
has been reduced to solving equations~\eqref{eq:7.1.24} and \eqref{eq:7.1.25} for
the two axisymmetric functions \(V\) and \(w\) in ordinary three-dimensional
Euclidean space. This is a remarkable simplification of the original problem of
solving the full set of Einstein's equation for 10 unknown functions \(g_{\mu\nu}\).
Nevertheless, these equations are still sufficiently difficult to solve that---with the
exception of static solutions discussed below---almost no solutions have been found
by direct attack on equations~\eqref{eq:7.1.24} and \eqref{eq:7.1.25}. (In cases
where solutions have been obtained, it often has proven useful to work with
coordinates other than \(\rho\) and \(z\) [Zipoy 1966; Chandrasekhar 1983].) The
equations can be reformulated through the introduction of potentials (Ernst 1968),
and a few solutions of interest have been found by direct study of these modified
equations (Tomimatsu and Sato 1972, 1973). However, recent progress in methods
for generating solutions (discussed briefly in section 7.4 below) has produced
algorithms for obtaining \emph{all} the asymptotically flat stationary, axisymmetric
solutions, although the computations required in this procedure remain formidable.

The major exception where direct attack has been completely successful is the case
of static, axisymmetric spacetimes (with the axisymmetric Killing field \(\psi^a\)
lying in the hypersurfaces orthogonal to the static Killing field). In that case, we
have \(w=0\), so equation~\eqref{eq:7.1.25} is trivially satisfied, and if we define
\(\chi=\ln V\), equation~\eqref{eq:7.1.24} reduces to simply
\begin{equation}
  \boldsymbol{\nabla}^2\chi=0 ,
  \tag{7.1.29}\label{eq:7.1.29}
\end{equation}
i.e., \(\chi\) is an axisymmetric solution of the ordinary Laplace equation in
three-dimensional flat space. Since all such solutions of this equation are explicitly
known, all static, axisymmetric solutions of Einstein's equation can be obtained.
This analysis of the static, axisymmetric spacetimes was first carried out by Weyl
(1917), and the solutions are often referred to as the Weyl solutions. It should be
pointed out that the properties of the solution of equation~\eqref{eq:7.1.29} do not
translate in a simple way into properties of the spacetime metric it generates.
Specifically, the monopole solution of Laplace's equation does \emph{not} generate
a spherically symmetric spacetime. To obtain the spherically symmetric
Schwarzschild solution by this procedure, one must choose \(\chi\) to be the
potential of a finite rod on the \(z\)-axis (centered on the origin) with constant mass
per unit length.
